% !TeX program = xelatex
% !TeX encoding = UTF-8

% --- 使用 ctexart 文档类，自动支持中文 ---
\documentclass[a4paper]{ctexart}

% --- 导入所需宏包 ---
\usepackage[a4paper, top=2.5cm, bottom=2.5cm, left=2cm, right=2cm]{geometry} % 页面边距
\usepackage{graphicx}   % 插入图片
\usepackage{booktabs}   % 专业表格
\usepackage{amsmath}    % 数学公式
\usepackage{amssymb}    % 更多数学符号 (例如 \nrightarrow)
\usepackage{listings}   % 排版代码
\usepackage{xcolor}     % 代码高亮
\usepackage{hyperref}   % PDF链接和书签

% --- PDF 书签设置 ---
\hypersetup{
	colorlinks=true,
	linkcolor=blue,
	urlcolor=blue,
	citecolor=blue,
	pdftitle={基于时间序列模型的比特币与以太坊价格关系研究},
	pdfauthor={林鑫},
	bookmarks=true,
	bookmarksopen=true
}

% --- Python 代码高亮样式定义 ---
\definecolor{codegreen}{rgb}{0,0.5,0}
\definecolor{codegray}{rgb}{0.5,0.5,0.5}
\definecolor{codepurple}{rgb}{0.58,0,0.82}
\definecolor{backcolour}{rgb}{0.98,0.98,0.98}

\lstdefinestyle{pythonstyle}{
	language=Python,
	backgroundcolor=\color{backcolour},   
	commentstyle=\color{codegreen},
	keywordstyle=\color{blue},
	stringstyle=\color{codepurple},
	basicstyle=\ttfamily\small,
	breakatwhitespace=false,         
	breaklines=true,                 
	captionpos=b,                    
	keepspaces=true,                 
	numbers=left,                    
	numbersep=5pt,                  
	showspaces=false,                
	showstringspaces=false,
	showtabs=false,                  
	tabsize=2,
	frame=single,
	framerule=0pt,
	rulesepcolor=\color{gray}
}
\lstset{style=pythonstyle} % 全局应用

% --- 封面信息 ---
\title{
	时间序列分析 \\
	课程论文 \\
	\vspace{1.5cm}
	\Huge 基于时间序列模型的比特币与以太坊价格关系研究
}
\author{
	班级：信计2302班 \\
	专业：信息与计算科学 \\
	学号：1043220122 \\
	姓名：林鑫
}
\date{2025年11月}

% --- 文档开始 ---
\begin{document}
	
	\maketitle
	\newpage
	
	\begin{abstract}
		\noindent 本研究基于 2020 年 1 月至 2025 年 10 月比特币（BTC）与以太坊（ETH）的日度收盘价数据，采用多元时间序列模型系统分析了其风险特征与动态关联。实证结果显示，两者的收益率序列均呈现显著的“尖峰厚尾”和“波动聚集”特征。针对此特性，本文构建了 ARIMA-GARCH-t 联合模型，有效捕捉了市场的条件异方差性，并实现了对 BTC 价格走势及 95\% 置信区间的有效预测。在资产关联性方面，Johansen 协整检验表明两者缺乏长期稳定的均衡关系；但 Granger 因果检验揭示了显著的短期双向信息溢出效应。进一步基于 VAR(10) 模型的脉冲响应与方差分解分析发现，市场存在显著的不对称传导机制：BTC 具有绝对主导权，其价格冲击能迅速同向传导至 ETH，解释了 ETH 约 67\% 的价格波动；而 ETH 对 BTC 的反向影响微弱。本研究量化揭示了加密市场“BTC引导、ETH跟随”的微观结构，为数字资产的风险管理与投资组合策略提供了实证依据。
		
		\vspace{1cm}
		\noindent \textbf{关键词：} 比特币；以太坊；ARIMA-GARCH-t；VAR模型；波动溢出；非对称传导
	\end{abstract}
	
	\newpage
	\tableofcontents % 生成目录
	\newpage
	
	\section{引言}
	随着数字经济的蓬勃发展，加密货币市场已成为全球金融体系中不可忽视的新兴板块。作为市值最高、影响力最大的两种核心资产，比特币（BTC）和以太坊（ETH）的价格行为不仅反映了市场情绪的变迁，也蕴含了复杂的风险传递机制。然而，加密市场典型的高波动性、“尖峰厚尾”分布以及资产间复杂的动态关联，使得传统的线性金融模型在分析和预测时面临严峻挑战。
	
	为了更准确地刻画和理解这两种资产的价格动态，本研究采用了多层次的时间序列分析框架。首先，针对单一资产收益率普遍存在的波动聚集和非正态分布特征，传统的 GARCH 模型往往难以充分捕捉极端风险。因此，本研究引入了基于 t-分布的 GARCH 模型（GARCH-t），旨在更精确地描述收益率的尾部特征和时变波动率。
	
	其次，在探究资产间关系时，区分长期均衡与短期联动至关重要。本研究通过 Johansen 协整检验考察了两者是否存在长期稳定的价格锚定关系，并通过 VAR 模型框架下的格兰杰因果检验、脉冲响应函数及方差分解技术，深入解构了两者在短期内的信息传导方向和强度，试图回答“谁在引领市场”这一核心问题。
	
	本报告旨在通过严谨的实证分析，揭示 BTC 与 ETH 市场的内在运行规律，验证成熟金融计量模型在新兴加密市场中的有效性，并为投资者进行风险规避和策略制定提供量化参考。
	
	\section{数据基本统计分析}
	\subsection{数据来源}
	本实验数据均来源于CoinGecko网站下载，截取了2020-01-01至2025-10-30的每日收盘价，单位为USDT（大致相当于美元），以下为数据预览：
	
	\begin{verbatim}
		date,btc_close,eth_close
		2020-01-01 00:00:00 UTC,7195.153895430029,129.18638529253852
		2020-01-02 00:00:00 UTC,7193.7546679601,130.4846850311141
		2020-01-03 00:00:00 UTC,6963.853353983485,127.04525801179804
		2020-01-04 00:00:00 UTC,7300.34693530511,133.70264861844188
		2020-01-05 00:00:00 UTC,7351.377053166627,134.1368826978575
	\end{verbatim}
	
	\subsection{数据基本统计分析数据}
	\subsubsection{基本数据特征}
	下表展示了 \texttt{btc\_close} (比特币收盘价) 和 \texttt{eth\_close} (以太坊收盘价) 的核心描述性统计数据。
	
	\begin{table}[htbp]
		\centering
		\caption{描述性统计数据}
		\label{tab:desc_stats}
		\begin{tabular}{lrr}
			\toprule
			统计量 & btc\_close & eth\_close \\
			\midrule
			计数 (count) & 2130.000000 & 2130.000000 \\
			均值 (mean) & 45832.318269 & 2134.436787 \\
			标准差 (std) & 30839.402336 & 1183.327106 \\
			最小值 (min) & 5032.501351 & 110.597898 \\
			25\%分位数 (25\%) & 21428.063943 & 1384.076317 \\
			中位数 (50\%) & 38646.722846 & 2002.150320 \\
			75\%分位数 (75\%) & 62868.922659 & 3056.962393 \\
			最大值 (max) & 124773.508231 & 4829.225542 \\
			\bottomrule
		\end{tabular}
	\end{table}
	
	\subsubsection{数据可视化}
	
	\begin{figure}[htbp]
		\centering
		% --- 图像占位符 ---
		% 请将您的图表导出为 'figure1.png' 并与本 .tex 文件放在同一目录
		% 然后取消下一行的注释
		%\includegraphics[width=\textwidth]{figure1.png}
		% --- 占位符 ---
		\framebox{\parbox{0.9\textwidth}{\centering
				\vspace{8cm} \textbf{图像占位符 (图1)} \\
				\small\textit{这里应插入“原始价格时序图 (BTC vs ETH)”}
				\vspace{8cm}
		}}
		% --- 占位符结束 ---
		\caption{原始价格时序图 (BTC vs ETH)}
		\label{fig:price_timeseries}
	\end{figure}
	
	\begin{figure}[htbp]
		\centering
		% --- 图像占位符 ---
		% 请将您的图表导出为 'figure4.png' 并与本 .tex 文件放在同一目录
		% 然后取消下一行的注释
		%\includegraphics[width=\textwidth]{figure4.png}
		% --- 占位符 ---
		\framebox{\parbox{0.9\textwidth}{\centering
				\vspace{12cm} \textbf{图像占位符 (图4)} \\
				\small\textit{这里应插入“原始价格的自相关性 (非平稳性展示)”}
				\vspace{12cm}
		}}
		% --- 占位符结束 ---
		\caption{原始价格的自相关性 (非平稳性展示)}
		\label{fig:price_acf}
	\end{figure}
	
	\paragraph{原始价格时序图及其自相关性分析}
	该图表（图\ref{fig:price_timeseries}）展示了从2020年初到2025年中期的比特币（btc\_close）和以太坊（eth\_close）的每日收盘价时间序列。通过观察图\ref{fig:price_timeseries}及其自相关图（图\ref{fig:price_acf}），我们可以总结出以下几个显著的时序特征：
	
	\begin{enumerate}
		\item \textbf{显著的长期趋势（Trend）：} 两个时间序列（BTC 和 ETH）均表现出明显的长期上升趋势。
		\item \textbf{周期性波动（Cyclicality）：} 两个序列都具有非常明显的周期性特征，即牛市和熊市的交替。
		\item \textbf{波动率聚集性（Volatility Clustering）：} 这是一个非常关键的金融时序特征。在价格快速上涨或快速下跌的时期，价格的日内或日间摆动幅度显著增大。
		\item \textbf{非平稳性（Non-Stationarity）：}
		\begin{itemize}
			\item \textbf{均值非平稳：} 由于存在明显的长期趋势，序列的均值随时间而变化。
			\item \textbf{方差非平稳：} 由于存在波动率聚集，序列的方差（波动性）也随时间而变化。
			\item \textbf{自相关性证据 (图\ref{fig:price_acf})：} ACF 缓慢衰减 和 PACF 一阶截尾是随机游走过程（Random Walk, I(1)过程）的经典标志。
		\end{itemize}
		\item \textbf{高度相关性（High Correlation）：} BTC 和 ETH 的价格走势表现出极强的正相关性。
	\end{enumerate}
	
	\subsubsection{平稳性与白噪声检验}
	为了在后续步骤中建立有效的时间序列模型，我们必须使用平稳的数据序列。
	
	\paragraph{1. 原始价格平稳性检验 (ADF)}
	ADF检验的原假设（H0）是：序列存在单位根，即序列非平稳。
	\begin{itemize}
		\item \textbf{BTC 价格:} `ADF 统计量`: -0.3899, `p-value`: 0.9118。结论: 无法拒绝原假设。序列非平稳。
		\item \textbf{ETH 价格:} `ADF 统计量`: -1.8763, `p-value`: 0.3433。结论: 无法拒绝原假设。序列非平稳。
	\end{itemize}
	
	\subsubsection{数据处理}
	考虑到金融数据的特征，我们采用对数据取对数后做差分，即对数收益率（ $R_t = \ln(P_t) - \ln(P_{t-1})$ ）。
	
	\begin{figure}[htbp]
		\centering
		% --- 图像占位符 ---
		% 请将您的图表导出为 'figure2.png' 并与本 .tex 文件放在同一目录
		% 然后取消下一行的注释
		%\includegraphics[width=\textwidth]{figure2.png}
		% --- 占位符 ---
		\framebox{\parbox{0.9\textwidth}{\centering
				\vspace{8cm} \textbf{图像占位符 (图2)} \\
				\small\textit{这里应插入“对数收益率时序图 (BTC vs ETH)”}
				\vspace{8cm}
		}}
		% --- 占位符结束 ---
		\caption{对数收益率时序图 (BTC vs ETH)}
		\label{fig:return_timeseries}
	\end{figure}
	
	\paragraph{1. 平稳性检验 (ADF)}
	\begin{itemize}
		\item \textbf{BTC 对数收益率:} `p-value`: 9.3209e-28。结论: 强烈拒绝原假设。序列已平稳。
		\item \textbf{ETH 对数收益率:} `p-value`: 1.8476e-22。结论: 强烈拒绝原假设。序列已平稳。
	\end{itemize}
	
	\paragraph{2. 白噪声检验 (Ljung-Box)}
	Ljung-Box 检验的原假设（H0）是：序列是白噪声。
	\begin{itemize}
		\item \textbf{BTC 对数收益率:} `p-value`: 0.00677。结论: 强烈拒绝原假设。序列非白噪声。
		\item \textbf{ETH 对数收益率:} `p-value`: 0.000062。结论: 强烈拒绝原假设。序列非白噪声。
	\end{itemize}
	
	\paragraph{3. 总体结论}
	ADF 检验证明，对数收益率序列是平稳的。Ljung-Box 检验证明，这个平稳的收益率序列并不是纯粹的随机白噪声。这一结果证实了我们的数据是“可建模”的。
	
	\subsubsection{检验尖峰厚尾与Q-Q图}
	\paragraph{1. 收益率分布特征分析 (尖峰厚尾)}
	我们必须检验这个收益率序列本身的分布特征。
	
	\begin{figure}[htbp]
		\centering
		% --- 图像占位符 ---
		% 请将您的图表导出为 'figure3.png' 并与本 .tex 文件放在同一目录
		% 然后取消下一行的注释
		%\includegraphics[width=\textwidth]{figure3.png}
		% --- 占位符 ---
		\framebox{\parbox{0.9\textwidth}{\centering
				\vspace{6cm} \textbf{图像占位符 (图3)} \\
				\small\textit{这里应插入“对数收益率分布直方图”}
				\vspace{6cm}
		}}
		% --- 占位符结束 ---
		\caption{对数收益率分布直方图}
		\label{fig:return_hist}
	\end{figure}
	
	\paragraph{统计检验：偏度 (Skewness) 与 峰度 (Kurtosis)}
	\begin{itemize}
		\item \textbf{偏度 (Skewness):} BTC: -1.1759, ETH: -1.1491。显著的负偏度（“左偏”），表明极端的负收益（暴跌）更为常见或幅度更大。
		\item \textbf{峰度 (Kurtosis):} BTC: 18.0235, ETH: 16.5673。极高的正峰度值，这是“尖峰厚尾”（Leptokurtosis / Fat Tails）的决定性统计证据。
	\end{itemize}
	
	\begin{figure}[htbp]
		\centering
		% --- 图像占位符 ---
		% 请将您的图表导出为 'figure3_1.png' 并与本 .tex 文件放在同一目录
		% 然后取消下一行的注释
		%\includegraphics[width=\textwidth]{figure3_1.png}
		% --- 占位符 ---
		\framebox{\parbox{0.9\textwidth}{\centering
				\vspace{6cm} \textbf{图像占位符 (图3.1)} \\
				\small\textit{这里应插入“Q-Q 图 (证明尖峰厚尾特性)”}
				\vspace{6cm}
		}}
		% --- 占位符结束 ---
		\caption{Q-Q 图 (证明尖峰厚尾特性)}
		\label{fig:qq_plot}
	\end{figure}
	
	\paragraph{图3.1：Q-Q 图}
	Q-Q图（图\ref{fig:qq_plot}）是检验“厚尾”特征的最直观工具。两个图（BTC 和 ETH）均呈现出清晰的 "S" 形，数据点系统性地偏离了红线，证明了厚尾特性。
	
	\paragraph{结论}
	对数收益率序列不服从正态分布，均表现出显著的负偏度和极其强烈的“尖峰厚尾”特性。任何假设正态分布的简单模型都将严重低估市场极端风险。
	
	\section{一元时间序列建模}
	\subsection{基本原理}
	\subsubsection{ARIMA模型原理}
	ARIMA(p, d, q)（自回归积分滑动平均）模型是经典的线性预测模型。其一般表达式为：
	$$ \phi(B)(1-B)^d y_t = \theta(B)\epsilon_t $$
	其中，$B$ 为滞后算子（ $By_t = y_{t-1}$ ）；$\phi(B)$ 与 $\theta(B)$ 分别为自回归与滑动平均多项式；$\epsilon_t$ 为白噪声项。
	
	\subsubsection{GARCH模型原理}
	GARCH (Generalized ARCH) 模型用于动态刻画条件方差的演化过程。
	均值方程：$ y_t = \mu + \epsilon_t $
	条件方差方程 (GARCH(1,1)):
	$$ \sigma_t^2 = \omega + \alpha \epsilon_{t-1}^2 + \beta \sigma_{t-1}^2 $$
	其中 $\alpha$ 代表短期冲击效应，$\beta$ 代表波动的持续效应。
	
	\subsubsection{ARIMA-GARCH混合模型}
	ARIMA模型用于刻画均值动态结构，GARCH模型用于刻画残差的方差动态结构。
	$$ y_t = \mu + \phi(B)(y_{t-1}) + \epsilon_t $$
	$$ \epsilon_t = \sigma_t z_t, \quad z_t \sim N(0,1) \text{ 或 } t(\nu) $$
	$$ \sigma_t^2 = \omega + \alpha \epsilon_{t-1}^2 + \beta \sigma_{t-1}^2 $$
	
	\subsection{ARIMA-GARCH建模}
	我们采用了 ARIMA-GARCH 的两步建模法。
	
	\subsubsection{1. 均值方程：ARIMA(1, 0, 1)}
	使用 \texttt{auto\_arima\_select} 函数，发现 ARMA(1, 1) 是最优模型。
	
	% --- 占位符：ARIMA 摘要 ---
	\begin{verbatim}
		--- BTC 对数收益率 ARIMA 模型摘要 ---
		==============================================================================
		Dep. Variable:        btc_log_return   No. Observations:                 2129
		Model:                 ARIMA(1, 0, 1)   Log Likelihood                4291.187
		...
		==============================================================================
		coef    std err          z      P>|z|      [0.025      0.975]
		------------------------------------------------------------------------------
		const          0.0013      0.001      1.808      0.071      -0.000       0.003
		ar.L1         -0.5954      0.185     -3.212      0.001      -0.959      -0.232
		ma.L1          0.5436      0.194      2.799      0.005       0.163       0.924
		sigma2         0.0010   1.07e-05     97.110      0.000       0.001       0.001
		==============================================================================
		Ljung-Box (L1) (Q):                   0.10   Jarque-Bera (JB):             29457.02
		Prob(Q):                              0.75   Prob(JB):                         0.00
		Heteroskedasticity (H):               0.38   Skew:                            -1.20
		Prob(H) (two-sided):                  0.00   Kurtosis:                        21.06
		==============================================================================
	\end{verbatim}
	% --- 占位符结束 ---
	
	% --- 占位符：ARIMA 检验 ---
	\begin{verbatim}
		BTC 对数收益率 参数显著性 (p-values):
		ar.L1                p = 0.001317 -> 显著
		ma.L1                p = 0.005133 -> 显著
		
		BTC 对数收益率 残差 Ljung-Box 检验 (lags = [10, 20]):
		lag=10 -> lb_stat=14.5256, p=0.150338 -> 通过(白噪声)
		lag=20 -> lb_stat=26.2755, p=0.156905 -> 通过(白噪声)
	\end{verbatim}
	% --- 占位符结束 ---
	
	ARIMA(1,1) 模型的残差已经是白噪声。但 `Prob(H) = 0.00` 和 `Prob(JB) = 0.00` 证明了残差中仍存在异方差和“尖峰厚尾”特性。
	
	\begin{figure}[htbp]
		\centering
		% --- 图像占位符 ---
		% 请将您的图表导出为 'arima_resid_qq.png' 并与本 .tex 文件放在同一目录
		% 然后取消下一行的注释
		%\includegraphics[width=0.6\textwidth]{arima_resid_qq.png}
		% --- 占位符 ---
		\framebox{\parbox{0.6\textwidth}{\centering
				\vspace{6cm} \textbf{图像占位符} \\
				\small\textit{这里应插入“BTC 对数收益率 残差 Q-Q 图”}
				\vspace{6cm}
		}}
		% --- 占位符结束 ---
		\caption{BTC 对数收益率 ARIMA 残差 Q-Q 图}
		\label{fig:arima_resid_qq}
	\end{figure}
	
	\subsubsection{2. 波动率方程：GARCH(1, 1) 与 t-分布}
	我们对 ARIMA 的残差拟合 GARCH(1,1)-t 模型。
	
	% --- 占位符：GARCH 摘要 ---
	\begin{verbatim}
		--- BTC GARCH(1,1)-t 模型摘要: ---
		Distribution:                Standardized Student's t
		==============================================================================
		coef    std err          t      P>|t|      95.0% Conf. Int.
		------------------------------------------------------------------------------
		Volatility Model
		------------------------------------------------------------------------------
		omega          0.1559  8.572e-02      1.819  6.891e-02 [-1.208e-02,  0.324]
		alpha[1]       0.0697  1.827e-02      3.815  1.363e-04 [3.388e-02,  0.105]
		beta[1]        0.9303  1.845e-02     50.428      0.000   [  0.894,  0.966]
		Distribution
		------------------------------------------------------------------------------
		nu             3.0730      0.205     14.981  9.724e-51   [  2.671,  3.475]
		==============================================================================
	\end{verbatim}
	% --- 占位符结束 ---
	
	\begin{itemize}
		\item \textbf{波动率聚集被捕捉：} `alpha[1]` (p<0.001) 和 `beta[1]` (p=0.000) 均高度显著。
		\item \textbf{尖峰厚尾被捕捉：} `nu` (t-分布自由度) 系数为 3.0730 且 p-value 极小，强力证明了使用 t-分布的必要性。
	\end{itemize}
	
	\begin{figure}[htbp]
		\centering
		% --- 图像占位符 ---
		% 请将您的图表导出为 'figure5.png' 并与本 .tex 文件放在同一目录
		% 然后取消下一行的注释
		%\includegraphics[width=\textwidth]{figure5.png}
		% --- 占位符 ---
		\framebox{\parbox{0.9\textwidth}{\centering
				\vspace{6cm} \textbf{图像占位符 (图5)} \\
				\small\textit{这里应插入“BTC GARCH(1,1)-t 标准化残差”}
				\vspace{6cm}
		}}
		% --- 占位符结束 ---
		\caption{BTC GARCH(1,1)-t 标准化残差}
		\label{fig:garch_resid}
	\end{figure}
	
	\subsection{最终模型表达式}
	\subsubsection{ARIMA(1, 0, 1) 模型}
	均值方程 ($Y_t$ 为对数收益率):
	$$ Y_t = 0.0013 - 0.5954 \cdot Y_{t-1} + 0.5436 \cdot \epsilon_{t-1} + \epsilon_t $$
	残差方差 $\sigma^2$ 估计值为 0.0010。
	
	\subsubsection{GARCH(1,1)-t模型}
	波动率方程 ($a_t$ 为 ARIMA 残差):
	$$ \sigma_t^2 = 0.1559 + 0.0697 a_{t-1}^2 + 0.9303 \sigma_{t-1}^2 $$
	分布假设：标准化残差 $\epsilon_t \sim t(3.0730)$。
	
	\subsection{模型预测}
	
	\begin{table}[htbp]
		\centering
		\caption{未来 10 天价格预测与 95\% 置信区间}
		\label{tab:price_forecast}
		\begin{tabular}{lrrr}
			\toprule
			日期(Date) & 预测价格 & 95\%置信下限 & 95\%置信上限 \\
			\midrule
			2025-10-31 & 110,331.40 & 102,013.99 & 119,326.95 \\
			2025-11-01 & 110,386.28 & 98,736.08 & 123,411.13 \\
			2025-11-02 & 110,578.50 & 96,378.35 & 126,870.86 \\
			2025-11-03 & 110,689.13 & 94,353.94 & 129,852.37 \\
			2025-11-04 & 110,848.69 & 92,626.32 & 132,655.93 \\
			2025-11-05 & 110,979.36 & 91,052.98 & 135,266.51 \\
			2025-11-06 & 111,127.55 & 89,627.68 & 137,784.80 \\
			2025-11-07 & 111,265.58 & 88,298.75 & 140,206.17 \\
			2025-11-08 & 111,409.96 & 87,059.49 & 142,571.23 \\
			2025-11-09 & 111,550.84 & 85,887.69 & 144,882.1 \\
			\bottomrule
		\end{tabular}
	\end{table}
	
	\begin{figure}[htbp]
		\centering
		% --- 图像占位符 ---
		% 请将您的图表导出为 'figure5_4.png' 并与本 .tex 文件放在同一目录
		% 然后取消下一行的注释
		%\includegraphics[width=\textwidth]{figure5_4.png}
		% --- 占位符 ---
		\framebox{\parbox{0.9\textwidth}{\centering
				\vspace{8cm} \textbf{图像占位符 (图5.4)} \\
				\small\textit{这里应插入“BTC 价格全历史样本内拟合与未来 10 天预测”}
				\vspace{8cm}
		}}
		% --- 占位符结束 ---
		\caption{BTC 价格全历史样本内拟合与未来 10 天预测}
		\label{fig:forecast_full}
	\end{figure}
	
	\begin{figure}[htbp]
		\centering
		% --- 图像占位符 ---
		% 请将您的图表导出为 'figure5_5.png' 并与本 .tex 文件放在同一目录
		% 然后取消下一行的注释
		%\includegraphics[width=\textwidth]{figure5_5.png}
		% --- 占位符 ---
		\framebox{\parbox{0.9\textwidth}{\centering
				\vspace{8cm} \textbf{图像占位符 (图5.5)} \\
				\small\textit{这里应插入“局部样本内拟合...与未来预测”}
				\vspace{8cm}
		}}
		% --- 占位符结束 ---
		\caption{局部样本内拟合 (带历史置信区间) 与未来预测}
		\label{fig:forecast_zoom}
	\end{figure}
	
	\subsection{风险价值 (VaR) 回测分析 (基于 GARCH(1,1)-t 模型)}
	\subsubsection{1. VaR 的概念与计算方法}
	风险价值 (Value at Risk, VaR) 是一种广泛使用的市场风险度量指标。本报告使用 GARCH(1,1)-t 模型来计算动态 VaR。
	计算公式为：
	$$ \text{VaR}_t^\alpha = \mu_t + \sigma_t \cdot F^{-1}(\alpha, \nu) $$
	其中：$\alpha = 5\%$；$\mu_t$ 是条件均值；$\sigma_t$ 是 GARCH 预测的条件波动率；$F^{-1}(\alpha, \nu)$ 是自由度为 $\nu$ 的 t-分布的 5\% 分位点。
	
	\subsubsection{2. 回测结果分析 (图 8)}
	\begin{figure}[htbp]
		\centering
		% --- 图像占位符 ---
		% 请将您的图表导出为 'figure8.png' 并与本 .tex 文件放在同一目录
		% 然后取消下一行的注释
		%\includegraphics[width=\textwidth]{figure8.png}
		% --- 占位符 ---
		\framebox{\parbox{0.9\textwidth}{\centering
				\vspace{8cm} \textbf{图像占位符 (图8)} \\
				\small\textit{这里应插入“BTC 每日收益率 vs. 95\% 动态风险价值 (VaR)”}
				\vspace{8cm}
		}}
		% --- 占位符结束 ---
		\caption{BTC 每日收益率 vs. 95\% 动态风险价值 (VaR)}
		\label{fig:var_backtest}
	\end{figure}
	图 \ref{fig:var_backtest} 直观地展示了 GARCH-VaR 模型的表现。红线（VaR阈值）随着蓝线（实际收益率）的波动而剧烈起伏，证明 GARCH 模型成功捕捉到了市场的波动率聚集特性。
	
	\subsubsection{3. 定量回测统计}
	\begin{itemize}
		\item 总观测天数：2129 天
		\item 理论突破率：5.00%
		\item 实际突破次数：32 次
		\item \textbf{实际突破率：} $32 \div 2129 \approx \mathbf{1.50\%}$
	\end{itemize}
	
	\paragraph{3.1 结果解读：模型过于保守}
	结果 (1.50\%) 显著低于理论值 (5.00%)。这意味着 GARCH-VaR 模型高估了市场风险。这可能因为模型拟合出的 t-分布自由度 $\nu \approx 3.07$ 非常低，过度夸大了极端黑天鹅事件的概率。
	
	\subsubsection{4. 未来一日风险预测}
	\begin{itemize}
		\item \textbf{95\% VaR: -5.7166\%}
		\item \textbf{解读：} 根据当前的市场波动率水平，模型有 95\% 的把握认为，明天 BTC 的跌幅不会超过 5.72%。
	\end{itemize}
	
	\section{多元时间序列建模}
	\subsection{协整性检验}
	\subsubsection{1. 滞后阶数选择}
	\begin{itemize}
		\item 收益率序列 (VAR): 最优滞后阶数为 10。
		\item 价格序列 (VAR): 最优滞后阶数为 8。
	\end{itemize}
	
	\subsubsection{2. 协整检验 (Engle-Granger 两步法)}
	% --- 占位符：协整检验 ---
	\begin{verbatim}
		--- 5.2 协整检验 (Engle-Granger 两步法) ---
		H0: 两个序列没有协整关系
		
		检验 1 (BTC ~ ETH):
		t-statistic: -1.1845
		p-value: 0.8619
		结论: 无法拒绝 H0 (p > 0.05)。序列无协整关系。
		---
		检验 2 (ETH ~ BTC):
		t-statistic: -2.3851
		p-value: 0.3319
		结论: 无法拒绝 H0 (p > 0.05)。序列无协整关系。
	\end{verbatim}
	% --- 占位符结束 ---
	\textbf{分析:} 协整检验结果表明，BTC 和 ETH 的价格在长期走势上并没有被绑定在一起。因此，建立向量误差修正模型 (VECM) 并不合适。
	
	\subsection{格兰杰因果检验 (Granger Causality Test)}
	我们在平稳的收益率序列上进行了格兰杰因果检验（滞后 10 阶）。
	\begin{itemize}
		\item \textbf{检验 1: ETH 是否格兰杰导致 BTC ($H_0$: ETH $\nrightarrow$ BTC)}
		\begin{itemize}
			\item `p-value` = 0.000003
			\item \textbf{结论:} 强烈拒绝原假设。ETH 的过去收益率对预测 BTC 有显著意义。
		\end{itemize}
		\item \textbf{检验 2: BTC 是否格兰杰导致 ETH ($H_0$: BTC $\nrightarrow$ ETH)}
		\begin{itemize}
			\item `p-value` = 0.008054
			\item \textbf{结论:} 拒绝原假设。BTC 的过去收益率对预测 ETH 有显著作用。
		\end{itemize}
	\end{itemize}
	\textbf{分析:} 格兰杰因果检验揭示了 BTC 和 ETH 市场之间存在显著的\textbf{双向因果关系}。
	
	\subsection{VAR（向量自回归）模型}
	\subsubsection{1. 模型选择依据}
	基于“无长期协整关系”和“存在短期因果关系”，向量自回归 (VAR) 模型成为最合适的选择。
	
	\subsubsection{2. 模型原理简述}
	一个 p 阶的 VAR模型可以表示为：
	$$ Y_t = C + A_1 Y_{t-1} + A_2 Y_{t-2} + \dots + A_p Y_{t-p} + \epsilon_t $$
	其中：$Y_t = (BTC\_Ret_t, ETH\_Ret_t)'$ 是内生变量向量；$A_1, \dots, A_p$ 是系数矩阵；$\epsilon_t$ 是误差项向量。
	
	\subsubsection{3. 模型拟合结果分析}
	我们拟合了 VAR(10) 模型。
	% --- 占位符：VAR 摘要 ---
	\begin{verbatim}
		--- 5.4 拟合 VAR(10) 模型 (在收益率上) ---
		Summary of Regression Results
		...
		Results for equation btc_log_return
		==============================================================================
		coefficient       std. err.      t-stat       prob.
		------------------------------------------------------------------------------
		...
		L1.eth_log_return  -0.076008       0.028296      -2.686      0.007
		...
		==============================================================================
		Results for equation eth_log_return
		==============================================================================
		coefficient       std. err.      t-stat       prob.
		------------------------------------------------------------------------------
		...
		L5.btc_log_return   0.138810       0.051952       2.670      0.008
		L6.btc_log_return  -0.164368       0.051747      -3.176      0.002
		...
		==============================================================================
		Correlation matrix of residuals
		btc_log_return   eth_log_return
		btc_log_return         1.000000         0.822717
		eth_log_return         0.822717         1.000000
	\end{verbatim}
	% --- 占位符结束 ---
	
	\paragraph{3.(1) BTC 收益率方程 (btc\_log\_return)}
	BTC 当期收益率受 ETH 滞后 1, 4, 5, 6, 7, 10 期的显著影响。
	$$ R_{BTC,t} = 0.0012 + \dots - 0.076 R_{ETH,t-1} + \dots + \epsilon_{BTC,t} $$
	\paragraph{3.(2) ETH 收益率方程 (eth\_log\_return)}
	ETH 当期收益率受 BTC 滞后 5, 6 期的显著影响。
	$$ R_{ETH,t} = 0.0014 + \dots + 0.139 R_{BTC,t-5} - 0.164 R_{BTC,t-6} + \dots + \epsilon_{ETH,t} $$
	\paragraph{3.(3) 残差相关性分析}
	两个方程的误差项 $\epsilon_{BTC,t}$ 和 $\epsilon_{ETH,t}$ 并非独立，它们之间存在很强的同期相关性：
	$$ \text{Corr}(\epsilon_{BTC,t}, \epsilon_{ETH,t}) = 0.8227 $$
	这表明每天有大量新信息会同时推动两种资产的价格。
	
	\subsection{脉冲响应分析（IRF）}
	\subsubsection{1. 分析目的与原理}
	脉冲响应函数 (IRF) 模拟当一个变量受到冲击时，其他变量在未来的动态反应路径。我们采用正交化脉冲响应 (Orthogonalized IRF)，并通过 Cholesky 分解设定变量排序为 \textbf{[BTC收益率, ETH收益率]}。
	这一排序假设：冲击当天，BTC 可以影响 ETH，但 ETH 不能反向影响 BTC（只能从第二天开始）。
	
	\subsubsection{2. BTC 冲击对 ETH 的影响}
	\begin{figure}[htbp]
		\centering
		% --- 图像占位符 ---
		% 请将您的图表导出为 'figure6_btc_to_eth.png' 并与本 .tex 文件放在同一目录
		% 然后取消下一行的注释
		%\includegraphics[width=\textwidth]{figure6_btc_to_eth.png}
		% --- 占位符 ---
		\framebox{\parbox{0.9\textwidth}{\centering
				\vspace{8cm} \textbf{图像占位符 (图6, 上半部分)} \\
				\small\textit{这里应插入“Impulse responses (btc\_log\_return -> eth\_log\_return)”}
				\vspace{8cm}
		}}
		% --- 占位符结束 ---
		\caption{BTC 冲击对 ETH 的响应}
		\label{fig:irf_btc_to_eth}
	\end{figure}
	\begin{itemize}
		\item \textbf{当期反应 (第 0 天)：} 当 BTC 受到正冲击时，ETH 在当天 (Lag 0) 就表现出极强且显著的正向反应。
		\item \textbf{后续反应 (第 1-20 天)：} 冲击的影响衰减得非常快，在第 1 天迅速回落至 0 附近。
	\end{itemize}
	\textbf{结论：} BTC 对 ETH 的冲击是“脉冲式”的——来得猛烈，去得也快。
	
	\subsubsection{3. ETH 冲击对 BTC 的影响}
	\begin{figure}[htbp]
		\centering
		% --- 图像占位符 ---
		% 请将您的图表导出为 'figure6_eth_to_btc.png' 并与本 .tex 文件放在同一目录
		% 然后取消下一行的注释
		%\includegraphics[width=\textwidth]{figure6_eth_to_btc.png}
		% --- 占位符 ---
		\framebox{\parbox{0.9\textwidth}{\centering
				\vspace{8cm} \textbf{图像占位符 (图6, 下半部分)} \\
				\small\textit{这里应插入“Impulse responses (eth\_log\_return -> btc\_log\_return)”}
				\vspace{8cm}
		}}
		% --- 占位符结束 ---
		\caption{ETH 冲击对 BTC 的响应}
		\label{fig:irf_eth_to_btc}
	\end{figure}
	\begin{itemize}
		\item \textbf{当期反应 (第 0 天)：} 由于 Cholesky 排序的设定，当期影响为 0。
		\item \textbf{后续反应 (第 1-20 天)：} BTC 的响应呈现“震荡”模式，但 95\% 置信区间绝大多数时候都包含了 0，说明在统计上并不显著。
	\end{itemize}
	\textbf{结论：} ETH 对 BTC 的溢出效应要弱得多。
	
	\subsubsection{4. 总结}
	脉冲响应分析清晰地揭示了加密货币市场中的\textbf{不对称信息传导机制}：BTC 具有主导地位，ETH 具有跟随特征。
	
	\subsection{方差分解分析 (Variance Decomposition, FEVD)}
	\subsubsection{1. 分析目的与原理}
	方差分解量化了冲击的“重要性”。它告诉我们：一个币种价格波动的百分之多少是由它自己的“新消息”造成的，又有百分之多少是由另一个币种的“新消息”造成的。
	
	\subsubsection{2. 结果分析 (基于未来 10 天)}
	\begin{figure}[htbp]
		\centering
		% --- 图像占位符 ---
		% 请将您的图表导出为 'figure7.png' 并与本 .tex 文件放在同一目录
		% 然后取消下一行的注释
		%\includegraphics[width=\textwidth]{figure7.png}
		% --- 占位符 ---
		\framebox{\parbox{0.9\textwidth}{\centering
				\vspace{8cm} \textbf{图像占位符 (图7)} \\
				\small\textit{这里应插入“VAR(10) 方差分解 (未来 10 天)”}
				\vspace{8cm}
		}}
		% --- 占位符结束 ---
		\caption{VAR(10) 方差分解 (未来 10 天)}
		\label{fig:fevd}
	\end{figure}
	
	\begin{enumerate}
		\item \textbf{BTC 收益率的方差分解 (FEVD for btc\_log\_return)}
		BTC 表现出极强的外生性。到第 10 天（Lag 9），BTC 自身冲击对预测方差的贡献率依然高达 \textbf{98.08\%}，而 ETH 的贡献率仅为 1.92\%。
		\item \textbf{ETH 收益率的方差分解 (FEVD for eth\_log\_return)}
		ETH 表现出很强的内生性。在第 0 天（当期），ETH 的价格波动中就有高达 \textbf{67.69\%} 是由 BTC 的同期冲击解释的，只有 32.31\% 是由 ETH 自身的特有冲击解释的。
	\end{enumerate}
	
	\subsubsection{3. 总结}
	方差分解的结果提供了强有力的证据，支持了“BTC 是市场领导者，ETH 是市场跟随者”这一假说。
	
	\subsection{VAR模型预测}
	\subsubsection{1. 分析目的与原理}
	本节对比 ARIMA-GARCH-t 模型（时变波动率）与标准 VAR 模型（恒定波动率）在风险评估上的优劣。
	
	\subsubsection{2. 结果分析}
	\begin{figure}[htbp]
		\centering
		% --- 图像占位符 ---
		% 请将您的图表导出为 'figure5_7.png' 并与本 .tex 文件放在同一目录
		% 然后取消下一行的注释
		%\includegraphics[width=\textwidth]{figure5_7.png}
		% --- 占位符 ---
		\framebox{\parbox{0.9\textwidth}{\centering
				\vspace{8cm} \textbf{图像占位符 (图5.7)} \\
				\small\textit{这里应插入“局部拟合与预测对比 (GARCH vs. VAR)”}
				\vspace{8cm}
		}}
		% --- 占位符结束 ---
		\caption{局部拟合与预测对比 (GARCH vs. VAR)}
		\label{fig:forecast_compare}
	\end{figure}
	
	图 \ref{fig:forecast_compare} 对比了两种模型的拟合与预测：
	\begin{itemize}
		\item \textbf{历史拟合置信区间 (关键差异)：} GARCH (浅绿色阴影) 是一个动态变化的区间，能随市场波动收缩或扩张，完美包裹了真实价格。而 VAR (浅橙色阴影) 是一个宽度恒定的区间，无法适应市场动态。
		\item \textbf{未来预测置信区间：} GARCH (浅红色阴影) 根据当前的高波动状态，发出了更强烈的风险预警（更宽的置信区间）。而 VAR (浅橙色阴影) 因其“恒定波动率”假设，显著低估了即期风险。
	\end{itemize}
	
	\subsubsection{3. 总结}
	GARCH 在风险刻画上完胜 VAR。ARIMA-GARCH 模型提供了动态的、更符合现实的风险区间（区间预测），是比标准 VAR 模型更优越的金融分析工具。
	
	\section{总结与展望}
	\subsection{主要研究结论}
	\begin{enumerate}
		\item \textbf{市场风险特征显著：} BTC 与 ETH 收益率均呈现“尖峰厚尾”和“波动聚集”特征。
		\item \textbf{缺乏长期价格锚定：} 协整检验表明两者不存在长期稳定的均衡关系。
		\item \textbf{短期内 BTC 具有绝对主导权：} 格兰杰检验、脉冲响应和方差分解一致确认了市场存在显著的\textbf{不对称信息传导}。BTC 是“价格发现者”，ETH 是“跟随者”。
	\end{enumerate}
	
	\subsection{投资启示与建议}
	\begin{itemize}
		\item 风险管理需重视极端尾部风险，建议采用基于 GARCH-t 模型的动态 VaR。
		\item 以 BTC 为核心的市场监测指标，BTC 的动向是整个板块的变盘前兆。
		\item 跨品种套利需谨慎，基于均值回归的配对交易策略可能面临长期无法回归的风险。
	\end{itemize}
	
	\subsection{研究局限性与未来展望}
	\begin{enumerate}
		\item \textbf{极端风险刻画的保守性：} 95\% VaR 的实际突破率 (1.50\%) 远低于理论值 (5.00%)，表明 GARCH-t 模型过于保守。未来可尝试 GED、偏 t 分布或极值理论 (EVT) 来更精准建模尾部。
		\item \textbf{高频数据的应用缺失：} 本研究仅使用日度数据。未来可引入小时级、分钟级数据，利用已实现波动率（Realized Volatility）等指标进行更精细的考察。
		\item \textbf{动态相关性的简化处理：} VAR 假设相关性恒定。未来可采用 DCC-GARCH 或 Copula 模型探究相关结构的时变特征。
		\item \textbf{外部宏观因素的未纳入：} VAR 残差高达 0.82 的同期相关性暗示了大量共同外部冲击的存在。未来应将美联储货币政策、通胀数据等宏观变量作为外生因子纳入模型。
	\end{enumerate}
	
	% --- 附录 ---
	\appendix
	\section{Python 代码}
	
	\begin{lstlisting}[language=Python, caption={Python 完整分析代码}, label={lst:pythoncode}]
		import pandas as pd
		import numpy as np
		import matplotlib.pyplot as plt
		import seaborn as sns
		import statsmodels.api as sm
		import itertools
		from statsmodels.tsa.stattools import adfuller, grangercausalitytests, kpss
		from statsmodels.tsa.arima.model import ARIMA
		from statsmodels.graphics.tsaplots import plot_acf, plot_pacf
		from statsmodels.tsa.api import VAR
		from statsmodels.tsa.vector_ar.vecm import VECM, coint_johansen
		from statsmodels.stats.diagnostic import acorr_ljungbox
		from arch import arch_model
		from scipy.stats import t
		import warnings
		
		warnings.filterwarnings("ignore")
		plt.rcParams['font.sans-serif'] = ['SimHei'] # 用来正常显示中文标签
		plt.rcParams['axes.unicode_minus'] = False # 用来正常显示负号
		
		# --- 1. 数据导入和预处理 ---
		file_name = r"C:\study\时间序列分析\结课作业\data.xlsx"
		df = pd.read_excel(file_name)
		df['date'] = pd.to_datetime(df['date'])
		df.set_index('date', inplace=True)
		df.sort_index(inplace=True)
		df_prices = df[['btc_close', 'eth_close']].copy().dropna()
		print("数据加载成功，预处理完成。")
		print("数据概览:")
		print(df_prices.head())
		print("\n数据信息:")
		df_prices.info()
		
		# --- 2. 数据基本统计分析 ---
		print("\n--- 2. 描述性统计 ---")
		print(df_prices.describe())
		
		# 绘制价格时序图
		print("正在显示图1: 原始价格时序图 (BTC vs ETH)...")
		df_prices.plot(subplots=True, figsize=(10, 8))
		plt.suptitle('图1: 原始价格时序图 (BTC vs ETH)')
		plt.show()
		
		# 计算对数收益率
		df_returns = np.log(df_prices).diff().dropna()
		df_returns.columns = ['btc_log_return', 'eth_log_return']
		print("正在显示图2: 对数收益率时序图 (BTC vs ETH)...")
		df_returns.plot(subplots=True, figsize=(10, 8))
		plt.suptitle('图2: 对数收益率时序图 (BTC vs ETH)')
		plt.show()
		
		# 绘制对数收益率直方图
		print("正在显示图3: 对数收益率分布直方图...")
		plt.figure(figsize=(12, 5))
		plt.subplot(1, 2, 1)
		sns.histplot(df_returns['btc_log_return'], kde=True, bins=30)
		plt.title('BTC 对数收益率分布')
		plt.subplot(1, 2, 2)
		sns.histplot(df_returns['eth_log_return'], kde=True, bins=30)
		plt.title('ETH 对数收益率分布')
		plt.suptitle('图3: 对数收益率分布直方图')
		plt.show()
		
		# --- 改进 1: 检验尖峰厚尾 (Skewness, Kurtosis) 和 Q-Q 图 ---
		print("\n--- 2.1 收益率分布检验 (尖峰厚尾) ---")
		print(f"BTC 收益率 偏度 (Skewness): {df_returns['btc_log_return'].skew():.4f} (正态=0)")
		print(f"BTC 收益率 峰度 (Kurtosis): {df_returns['btc_log_return'].kurtosis():.4f} (正态=0, 金融>0)")
		print(f"ETH 收益率 偏度 (Skewness): {df_returns['eth_log_return'].skew():.4f}")
		print(f"ETH 收益率 峰度 (Kurtosis): {df_returns['eth_log_return'].kurtosis():.4f}")
		
		print("\n正在显示图3.1: 收益率 Q-Q 图 (对比正态分布)...")
		fig, axes = plt.subplots(1, 2, figsize=(12, 5))
		sm.qqplot(df_returns['btc_log_return'].dropna(), line='s', ax=axes[0])
		axes[0].set_title('BTC 收益率 Q-Q 图')
		sm.qqplot(df_returns['eth_log_return'].dropna(), line='s', ax=axes[1])
		axes[1].set_title('ETH 收益率 Q-Q 图')
		plt.suptitle('图3.1: Q-Q 图 (证明尖峰厚尾特性)')
		plt.show()
		
		# --- 3. 平稳性检验和白噪声检验 ---
		print("\n--- 3. 平稳性检验 (ADF) 与白噪声检验 ---")
		
		# --- 检验函数定义 ---
		# ADF 检验
		def run_adf_test(series, name):
		"""
		ADF 检验
		H0 (原假设): 序列存在单位根 (非平稳)
		"""
		result = adfuller(series)
		print(f'序列 "{name}" 的 ADF 检验 (H0: 存在单位根):')
		print(f'  ADF 统计量: {result[0]:.4f}, p-value: {result[1]:.4e}')
		if result[1] > 0.05:
		print(f'  结论: 序列非平稳 (p > 0.05)')
		else:
		print(f'  结论: 序列平稳 (p <= 0.05)')
		return result[1]
		
		# 白噪声检验 (Ljung-Box)
		def run_ljung_box_test(series, name, lags=10):
		"""
		Ljung-Box 白噪声检验
		H0 (原假设): 序列是白噪声 (即序列值彼此独立, 无自相关)
		"""
		# Ljung-Box test
		result_df = acorr_ljungbox(series, lags=[lags], return_df=True)
		p_value = result_df.iloc[0]['lb_pvalue']
		print(f'序列 "{name}" 的 Ljung-Box 白噪声检验 (H0: 序列为白噪声):')
		print(f'  L-B 统计量 (lag={lags}): {result_df.iloc[0]["lb_stat"]:.4f}, p-value: {p_value:.4e}')
		
		if p_value <= 0.05:
		print(f'  结论: 序列存在自相关 (非白噪声, 可建模) (p <= 0.05)')
		else:
		print(f'  结论: 序列为白噪声 (无信息) (p > 0.05)')
		return p_value
		
		# --- 3.1 原始价格检验 ---
		print("\n--- 3.1 原始价格平稳性检验 (预期: 非平稳) ---")
		p_btc_price_adf = run_adf_test(df_prices['btc_close'], 'BTC 价格')
		print("---")
		p_eth_price_adf = run_adf_test(df_prices['eth_close'], 'ETH 价格')
		
		# --- 3.2 对数收益率检验 ---
		print("\n--- 3.2 对数收益率平稳性与白噪声检验 (预期: 平稳, 非白噪声) ---")
		p_btc_return_adf = run_adf_test(df_returns['btc_log_return'], 'BTC 对数收益率')
		p_btc_return_lb = run_ljung_box_test(df_returns['btc_log_return'], 'BTC 对数收益率')
		print("---")
		p_eth_return_adf = run_adf_test(df_returns['eth_log_return'], 'ETH 对数收益率')
		p_eth_return_lb = run_ljung_box_test(df_returns['eth_log_return'], 'ETH 对数收益率')
		
		# --- 4. 一元时间序列模型 (ARIMA & GARCH) ---
		print("\n--- 4. 一元时间序列模型 (ARIMA & 升级版 GARCH) ---")
		if True:
		# --- 自动寻优的函数 (遍历 p,q) ---
		def auto_arima_select(ts, d=0, max_p=3, max_q=3):
		best_aic = np.inf
		best_bic = np.inf
		best_aic_res = None
		best_bic_res = None
		best_aic_order = None
		best_bic_order = None
		
		for p in range(max_p + 1):
		for q in range(max_q + 1):
		# 增加 (0,0,0) 的跳过，因为 GARCH 会处理 const
		if p == 0 and q == 0:
		continue
		try:
		model = ARIMA(ts, order=(p, d, q))
		res = model.fit()
		if res.aic < best_aic:
		best_aic = res.aic
		best_aic_res = res
		best_aic_order = (p, d, q)
		if res.bic < best_bic:
		best_bic = res.bic
		best_bic_res = res
		best_bic_order = (p, d, q)
		except Exception as e:
		continue
		return {
			'best_aic_order': best_aic_order,
			'best_bic_order': best_bic_order,
			'best_aic_result': best_aic_res,
			'best_bic_result': best_bic_res
		}
		
		# --- 残差诊断与显著性检验函数 ---
		def arima_diagnostics(result, series_name="Series", lb_lags=[10, 20]):
		print(f"\n--- {series_name} ARIMA 模型摘要 ---")
		print(result.summary())
		
		print(f"\n{series_name} 参数显著性 (p-values):")
		pvals = result.pvalues
		for name, pval in pvals.items():
		print(f"  {name:20s} p = {pval:.6f} -> {'显著' if pval < 0.05 else '不显著'}")
		
		print(f"\n{series_name} 残差 Ljung-Box 检验 (lags = {lb_lags}):")
		for lag in lb_lags:
		try:
		# 使用 statsmodels 0.13+ 的 acorr_ljungbox 格式
		lb = acorr_ljungbox(result.resid, lags=[lag], return_df=True)
		pvalue = lb['lb_pvalue'].iloc[0]
		print(f"  lag={lag} -> lb_stat={lb['lb_stat'].iloc[0]:.4f}, p={pvalue:.6f} -> {'通过(白噪声)' if pvalue>0.05 else '未通过(非白噪声)'}")
		except Exception as e:
		print(f"  lag={lag} 检验失败: {e}")
		
		# 残差绘图
		resid = result.resid
		fig, axes = plt.subplots(2, 2, figsize=(12, 8))
		axes[0, 0].plot(resid); axes[0, 0].set_title(f"{series_name} 残差时序")
		plot_acf(resid, lags=30, ax=axes[0, 1]); axes[0, 1].set_title(f"{series_name} 残差 ACF")
		plot_pacf(resid, lags=30, ax=axes[1, 0]); axes[1, 0].set_title(f"{series_name} 残差 PACF")
		sm.qqplot(resid, line='s', ax=axes[1, 1]); axes[1, 1].set_title(f"{series_name} 残差 Q-Q 图")
		plt.suptitle(f"{series_name} ARIMA 残差诊断")
		plt.tight_layout(rect=[0, 0.03, 1, 0.95])
		plt.show()
		
		# ========== 4.1 BTC ARIMA ==========
		print("\n--- 4.1 BTC 对数收益率 ARIMA 自动寻优 (基于 AIC/BIC) ---")
		# 确保 df_returns 存在 (在 Section 2 中创建)
		btc_res = auto_arima_select(df_returns['btc_log_return'], d=0, max_p=3, max_q=3)
		print(f"BTC 最优 AIC 阶数: {btc_res['best_aic_order']}")
		print(f"BTC 最优 BIC 阶数: {btc_res['best_bic_order']}")
		
		if btc_res['best_aic_result'] is not None:
		arima_result = btc_res['best_aic_result'] # 用于预测
		arima_diagnostics(arima_result, series_name="BTC 对数收益率")
		else:
		print("警告: 未能找到可拟合的 BTC ARIMA 模型。")
		
		# --- 4.3 GARCH(1,1)-t 模型 (BTC 对数收益率) ---
		print("\n--- 4.3 GARCH(1,1) 模型 (BTC 对数收益率) ---")
		try:
		garch_model_btc = arch_model(df_returns['btc_log_return'].dropna() * 100, 
		vol='Garch', p=1, o=0, q=1, dist='t')
		garch_result_btc = garch_model_btc.fit(update_freq=5, disp='off')
		
		print("\nBTC GARCH(1,1)-t 模型摘要:")
		print(garch_result_btc.summary())
		
		print("正在显示图5: BTC GARCH(1,1)-t 标准化残差...")
		garch_std_resid_btc = garch_result_btc.resid / garch_result_btc.conditional_volatility
		plt.figure(figsize=(10, 4))
		plt.plot(garch_std_resid_btc)
		plt.title('图5: BTC GARCH(1,1)-t 标准化残差')
		plt.show()
		
		except Exception as e:
		print(f"GARCH(1,1) 拟合失败: {e}")
		
		
		# --- 4.4 ARIMA-GARCH 联合模型 (样本内拟合与样本外预测) ---
		print("\n--- 4.4 ARIMA-GARCH 联合模型 (样本内拟合与样本外预测) ---")
		
		if 'arima_result' in locals() and 'garch_result_btc' in locals() and 'df_prices' in locals():
		try:
		last_price = df_prices['btc_close'].iloc[-1]
		last_date = df_prices.index[-1]
		forecast_horizon = 10
		
		# 从 GARCH 结果获取 t 分布自由度
		nu = garch_result_btc.params['nu']
		t_critical = t.ppf(0.975, nu)
		
		# --- 样本内拟合 ---
		in_sample_mu = arima_result.fittedvalues
		
		# 修正：样本内价格拟合应该基于 (t-1) 期的真实价格
		prev_actual_prices = df_prices['btc_close'].shift(1)
		in_sample_fitted_prices = prev_actual_prices * np.exp(in_sample_mu)
		
		# 样本内波动率 (GARCH 条件波动率)
		# 注意：GARCH 模型是用 100*return 拟合的，所以波动率要 /100
		in_sample_vol = garch_result_btc.conditional_volatility / 100
		
		# 样本内置信区间
		in_sample_upper = prev_actual_prices * np.exp(in_sample_mu + t_critical * in_sample_vol)
		in_sample_lower = prev_actual_prices * np.exp(in_sample_mu - t_critical * in_sample_vol)
		
		# --- 样本外预测 ---
		# 1. 均值预测 (来自 ARIMA)
		arima_forecast = arima_result.get_forecast(steps=forecast_horizon)
		pred_returns = arima_forecast.predicted_mean
		
		# 2. 波动率预测 (来自 GARCH)
		garch_forecast = garch_result_btc.forecast(horizon=forecast_horizon, reindex=False)
		pred_volatility = np.sqrt(garch_forecast.variance.iloc[-1].values) / 100
		
		# 3. 价格预测 (迭代)
		cum_pred_returns = np.cumsum(pred_returns)
		# 累计波动率 (注意：方差是可加的，波动率不是)
		cum_pred_volatility = np.sqrt(np.cumsum(pred_volatility**2))
		
		pred_prices = last_price * np.exp(cum_pred_returns)
		upper_prices = last_price * np.exp(cum_pred_returns + t_critical * cum_pred_volatility)
		lower_prices = last_price * np.exp(cum_pred_returns - t_critical * cum_pred_volatility)
		
		future_dates = pd.date_range(start=last_date + pd.Timedelta(days=1), periods=forecast_horizon, freq='D')
		
		price_forecast_df = pd.DataFrame({
			'Predicted_Price': pred_prices,
			'Lower_Price_95': lower_prices,
			'Upper_Price_95': upper_prices
		}, index=future_dates)
		
		print(f"\n未来 {forecast_horizon} 天预测结果:")
		print(price_forecast_df)
		
		# --- 绘图 1: 全历史 ---
		print(f"正在显示图5.4: 全历史样本内拟合与样本外预测...")
		plt.figure(figsize=(14, 7))
		plt.plot(df_prices.index, df_prices['btc_close'], color='blue', alpha=0.5, label='历史真实价格')
		plt.plot(in_sample_fitted_prices.index, in_sample_fitted_prices, color='green', linestyle='--', alpha=0.8, linewidth=1, label='样本内拟合价格')
		plt.plot(future_dates, price_forecast_df['Predicted_Price'], color='red', marker='o', markersize=4, label='未来预测价格')
		plt.fill_between(future_dates, price_forecast_df['Lower_Price_95'], price_forecast_df['Upper_Price_95'], color='red', alpha=0.15, label='未来 95% 置信区间')
		plt.title(f'图5.4: BTC 价格全历史样本内拟合与未来 {forecast_horizon} 天预测')
		plt.xlabel('日期')
		plt.ylabel('BTC 价格 (USD)')
		plt.legend(loc='upper left')
		plt.grid(True, alpha=0.3)
		plt.tight_layout()
		plt.show()
		
		# --- 绘图 2: 局部放大 (2025-09-01 至今) ---
		print(f"正在显示图5.5: 局部样本内拟合与样本外预测 (2025年9月至今)...")
		start_date_zoom = '2025-09-01'
		df_zoom = df_prices[start_date_zoom:]
		fitted_zoom = in_sample_fitted_prices[start_date_zoom:]
		lower_zoom = in_sample_lower[start_date_zoom:]
		upper_zoom = in_sample_upper[start_date_zoom:]
		
		if not df_zoom.empty:
		plt.figure(figsize=(14, 7))
		# 历史
		plt.plot(df_zoom.index, df_zoom['btc_close'], color='blue', marker='.', linestyle='-', linewidth=2, alpha=0.6, label='真实价格')
		plt.plot(fitted_zoom.index, fitted_zoom, color='green', linestyle='--', linewidth=2.5, alpha=0.9, label='拟合价格')
		plt.fill_between(fitted_zoom.index, lower_zoom, upper_zoom, color='green', alpha=0.1, label='历史 95% 置信区间')
		
		# 预测
		plt.plot(future_dates, price_forecast_df['Predicted_Price'], color='red', marker='o', markersize=6, linewidth=2, label='未来预测')
		plt.fill_between(future_dates, price_forecast_df['Lower_Price_95'], price_forecast_df['Upper_Price_95'], color='red', alpha=0.25, label='未来 95% 置信区间')
		
		plt.title(f'图5.5: 局部样本内拟合 (带历史置信区间) 与未来预测')
		plt.xlabel('日期')
		plt.ylabel('BTC 价格 (USD)')
		plt.legend(loc='upper left')
		plt.grid(True, alpha=0.3)
		plt.tight_layout()
		plt.xticks(rotation=45)
		plt.show()
		else:
		print(f"警告: 数据中没有 {start_date_zoom} 之后的数据。")
		
		except Exception as e:
		print(f"模型拟合与预测失败: {e}")
		else:
		print(f"错误: 缺少必要模型或数据。")
		
		
		# --- 5. 多元时间序列分析 (VECM/VAR) ---
		print("\n--- 5. 多元时间序列分析 (VECM / VAR) ---")
		if 'df_prices' in locals() and 'df_returns' in locals():
		
		# --- 5.1 确定最优滞后阶数 ---
		print("\n--- 5.1 确定最优滞后阶数 (Lag Selection) ---")
		
		# A. VAR(df_returns) 滞后阶数
		model_returns = VAR(df_returns)
		lag_selection_returns = model_returns.select_order(maxlags=10)
		optimal_lag_returns = lag_selection_returns.aic
		print(f"VAR(df_returns) 最优滞后 (AIC): {optimal_lag_returns}")
		
		# B. VAR(df_prices) 滞后阶数
		model_prices = VAR(df_prices)
		lag_selection_prices = model_prices.select_order(maxlags=10)
		optimal_lag_prices = lag_selection_prices.aic
		print(f"VAR(df_prices) 最优滞后 (AIC): {optimal_lag_prices}")
		
		
		# --- 5.2 协整检验 (Engle-Granger 两步法) ---
		print("\n--- 5.2 协整检验 (Engle-Granger 两步法) ---")
		print("H0: 两个序列没有协整关系")
		
		# EG 两步法 (statsmodels.tsa.stattools.coint)
		# 检验 1: btc_close = f(eth_close)
		print("\n检验 1 (BTC ~ ETH):")
		# trend='c' 代表包含常数项 (截距)
		# 返回值: (t-statistic, p-value, critical_values)
		score, pvalue, _ = sm.tsa.stattools.coint(df_prices['btc_close'], df_prices['eth_close'], trend='c')
		
		print(f'  t-statistic: {score:.4f}')
		print(f'  p-value: {pvalue:.4f}')
		if pvalue < 0.05:
		print(f'  结论: 拒绝 H0 (p < 0.05)。序列存在协整关系。')
		else:
		print(f'  结论: 无法拒绝 H0 (p > 0.05)。序列无协整关系。')
		
		print("---")
		
		# 检验 2: eth_close = f(btc_close)
		print("检验 2 (ETH ~ BTC):")
		score_rev, pvalue_rev, _ = sm.tsa.stattools.coint(df_prices['eth_close'], df_prices['btc_close'], trend='c')
		
		print(f'  t-statistic: {score_rev:.4f}')
		print(f'  p-value: {pvalue_rev:.4f}')
		if pvalue_rev < 0.05:
		print(f'  结论: 拒绝 H0 (p < 0.05)。序列存在协整关系。')
		else:
		print(f'  结论: 无法拒绝 H0 (p > 0.05)。序列无协整关系。')
		
		
		# --- 5.3 格兰杰因果检验 (Granger Causality Test) ---
		print("\n--- 5.3 格兰杰因果检验 (Granger Causality Test) ---")
		print(f"检验 'df_returns' (maxlag = {optimal_lag_returns})")
		
		# 检验 1: ETH 是否是 BTC 的格兰杰原因
		gc_result_eth_to_btc = grangercausalitytests(
		df_returns[['btc_log_return', 'eth_log_return']], 
		maxlag=optimal_lag_returns, 
		verbose=False
		)
		f_test_p_value_1 = gc_result_eth_to_btc[optimal_lag_returns][0]['ssr_ftest'][1]
		print(f"\nH0: ETH -/-> BTC (p-value): {f_test_p_value_1:.6f}")
		
		# 检验 2: BTC 是否是 ETH 的格兰杰原因
		gc_result_btc_to_eth = grangercausalitytests(
		df_returns[['eth_log_return', 'btc_log_return']], 
		maxlag=optimal_lag_returns, 
		verbose=False
		)
		f_test_p_value_2 = gc_result_btc_to_eth[optimal_lag_returns][0]['ssr_ftest'][1]
		print(f"\nH0: BTC -/-> ETH (p-value): {f_test_p_value_2:.6f}")
		
		# --- 5.4 拟合 VAR 模型 ---
		print("\n--- 5.4 拟合 VAR 模型 (在收益率上) ---")
		
		if 'df_returns' in locals() and 'optimal_lag_returns' in locals():
		try:
		# 1. 创建 VAR 模型实例
		model_var = VAR(df_returns)
		
		# 2. 拟合模型，使用之前确定的最优滞后阶数 (10)
		print(f"正在拟合 VAR({optimal_lag_returns}) 模型...")
		var_result = model_var.fit(optimal_lag_returns)
		
		# 3. 输出模型摘要
		print("\nVAR 模型摘要:")
		print(var_result.summary())
		
		# 4. 检查模型的稳定性 (根是否在单位圆内)
		print("\nVAR 模型稳定性检查 (所有根的模应 < 1):")
		is_stable = var_result.is_stable()
		print(f"模型是否稳定: {is_stable}")
		
		except Exception as e:
		print(f"VAR 模型拟合失败: {e}")
		else:
		print("错误: 缺少必要变量，无法拟合 VAR 模型。")
		
		
		# --- 5.5 脉冲响应分析 (IRF) ---
		print("\n--- 5.5 脉冲响应分析 (Impulse Response Analysis) ---")
		
		if 'var_result' in locals():
		try:
		# 设定分析的期数 (例如未来 20 天)
		irf_periods = 20
		
		# 'orth=True' 表示我们想要 Cholesky 分解后的正交化脉冲响应
		irf = var_result.irf(irf_periods)
		
		print(f"正在显示图6: 正交化脉冲响应函数 (未来 {irf_periods} 期)...")
		
		# 绘制所有组合
		irf.plot(orth=True, signif=0.05, figsize=(12, 10))
		plt.suptitle('图6: VAR(10) 正交化脉冲响应分析 (95% 置信区间)')
		plt.tight_layout(rect=[0, 0.03, 1, 0.95])
		plt.show()
		
		except Exception as e:
		print(f"脉冲响应分析失败: {e}")
		else:
		print("错误: 缺少 VAR 模型结果 (var_result)，无法进行 IRF 分析。")
		
		# --- 5.6 方差分解分析 (FEVD) ---
		print("\n--- 5.6 方差分解分析 (Variance Decomposition) ---")
		
		if 'var_result' in locals():
		try:
		# 设定分析期数 (例如未来 10 天)
		fevd_periods = 10
		
		# 计算方差分解
		fevd = var_result.fevd(fevd_periods)
		
		print(f"\n方差分解结果 (未来 {fevd_periods} 天):")
		print(fevd.summary())
		
		print(f"正在显示图7: 方差分解堆叠图...")
		# 绘制方差分解图
		fevd.plot(figsize=(12, 8))
		plt.suptitle(f'图7: VAR(10) 方差分解 (未来 {fevd_periods} 天)')
		plt.tight_layout(rect=[0, 0.03, 1, 0.95])
		plt.show()
		
		except Exception as e:
		print(f"方差分解分析失败: {e}")
		else:
		print("错误: 缺少 VAR 模型结果 (var_result)，无法进行 FEVD 分析。")
		
		else:
		print("错误: 缺少 'df_prices' 或 'df_returns' DataFrame，跳过多维分析。")
		
		
		# --- 6. 风险价值 (VaR) 计算与回测 ---
		print("\n--- 6. 风险价值 (VaR) 计算与回测 (基于 GARCH-t 模型) ---")
		
		if 'garch_result_btc' in locals():
		try:
		# 1. 获取模型参数
		garch_params = garch_result_btc.params
		mu = garch_params['mu']
		nu = garch_params['nu'] # t分布自由度
		
		# 2. 获取条件波动率 (Conditional Volatility)
		conditional_volatility = garch_result_btc.conditional_volatility
		
		# 3. 计算 95% VaR (alpha = 0.05)
		alpha = 0.05
		# t.ppf 是 t 分布的百分位点函数 (逆累积分布函数)
		q_t = t.ppf(alpha, nu)
		
		# VaR = E[R] + sigma_t * F_inv(alpha)
		# 注意: GARCH 模型是用 100*R 拟合的
		# 所以 mu 和 vol 都是基于 100*R 的尺度
		# VaR (in 100*R scale) = mu + conditional_volatility * q_t
		# VaR (in R scale) = (mu + conditional_volatility * q_t) / 100
		var_95 = (mu + conditional_volatility * q_t) / 100
		
		# 4. 创建用于绘图的 DataFrame
		var_df = pd.DataFrame({
			'Actual_Return': df_returns['btc_log_return'],
			'VaR_95': var_95
		}, index=df_returns.index).dropna()
		
		# 5. 绘制 VaR 回测图
		print("正在显示图8: BTC 每日收益率与 95% VaR 对比图...")
		plt.figure(figsize=(12, 6))
		plt.plot(var_df.index, var_df['Actual_Return'], color='blue', alpha=0.5, label='实际收益率')
		plt.plot(var_df.index, var_df['VaR_95'], color='red', linewidth=2, label='95% VaR (GARCH-t)')
		plt.title('图8: BTC 每日收益率 vs. 95% 动态风险价值 (VaR)')
		plt.legend()
		plt.grid(True, alpha=0.3)
		plt.show()
		
		# 6. 计算 VaR 突破次数 (Backtesting)
		exceptions = var_df[var_df['Actual_Return'] < var_df['VaR_95']]
		num_exceptions = len(exceptions)
		total_obs = len(var_df)
		actual_alpha = num_exceptions / total_obs
		
		print(f"\nVaR 回测结果 (95% 置信水平, 理论突破率应接近 5%):")
		print(f"  总观测天数: {total_obs}")
		print(f"  实际突破次数: {num_exceptions}")
		print(f"  实际突破率: {actual_alpha:.2%}")
		
		if abs(actual_alpha - 0.05) < 0.01: # 简单判断，实际应做 Kupiec 检验
		print("  -> 模型表现良好 (实际突破率接近理论值 5%)")
		else:
		print("  -> 模型可能需要改进 (实际突破率偏离理论值)")
		
		# 7. 预测下一交易日的 VaR
		forecast = garch_result_btc.forecast(horizon=1)
		next_vol = np.sqrt(forecast.variance.iloc[-1, 0])
		next_var_95 = (mu + next_vol * q_t) / 100
		
		print(f"\n未来一日风险预测:")
		print(f"  下一交易日 95% VaR: {next_var_95:.4%}")
		print(f"  (解读: 有 95% 的把握，明天 BTC 的损失不会超过 {abs(next_var_95):.2%})")
		
		except Exception as e:
		print(f"VaR 计算失败: {e}")
		else:
		print("错误: 缺少 GARCH 模型结果 (garch_result_btc)，无法计算 VaR。")
	\end{lstlisting}
	
\end{document}
% --- 文档结束 ---